\section{Заключение}

В работе спроектирован и реализован программный комплекс \textbf{Chronos} --- платформа для централизованного управления Docker Compose-проектами на нескольких узлах с веб-интерфейсом, динамической TCP-маршрутизацией через HAProxy, политиками резервного копирования томов в объектное хранилище и централизованным сбором логов в связке Loki/Grafana.

Достигнуты следующие результаты:
\begin{itemize}[leftmargin=1.5cm]
  \item обоснована актуальность подхода «compose + управляющий контур'' на фоне Kubernetes и классических панелей управления Docker;
  \item сформулированы и реализованы требования к модульной архитектуре (Core, Agent, Master, Web);
  \item реализованы публикация проектов, удалённое управление жизненным циклом, регистрация агентов, декларативные проверки и артефакты деплоя;
  \item внедрены механизмы снимков томов, фоновой репликации по политикам Master и интеграции с S3-совместимым API;
  \item реализован реестр TCP-маршрутов с генерацией конфигурации HAProxy и сценарием динамического применения на стенде;
  \item подтверждена работоспособность решения на локальном docker-compose стенде.
\end{itemize}

Практический итог --- работающий комплекс, который \textbf{уже применяется автором локально} для управления собственной инфраструктурой разработки: развёртывание сервисов, маршрутизация портов и наблюдаемость сведены к единому контуру без обязательного перехода на Kubernetes.

Направления развития включают усиление отказоустойчивости Master, расширение модели доступа, автоматизацию подсказок по портам из состояния compose и расширенные нагрузочные испытания.

\textbf{Цель работы достигнута}, поставленные задачи решены.
